\documentclass[11pt]{article}
\usepackage[a4paper,margin=1in]{geometry}
\usepackage{fancyhdr}
\usepackage{lastpage}
\usepackage{titlesec}
\usepackage{booktabs}
\usepackage{enumitem}
\usepackage{hyperref}
\usepackage{graphicx}
\usepackage{amsmath}
\usepackage{xcolor}
\usepackage{listings}

% Dark theme settings
\definecolor{darkbg}{rgb}{0.15,0.15,0.15}
\definecolor{darkfg}{rgb}{0.9,0.9,0.9}
\definecolor{keywordcolor}{rgb}{0,0.7,0.7} % cyan
\definecolor{stringcolor}{rgb}{0,0.7,0}   % green
\definecolor{commentcolor}{rgb}{0.5,0.5,0.5} % gray

\lstset{
    basicstyle=\ttfamily\color{darkfg},
    keywordstyle=\color{keywordcolor},
    stringstyle=\color{stringcolor},
    commentstyle=\color{commentcolor},
    backgroundcolor=\color{darkbg},
    frame=single,
    rulecolor=\color[rgb]{0.5,0.5,0.5},
    breaklines=true,
    showstringspaces=false
}

% Page color for dark theme
\usepackage{pagecolor}
\pagecolor{darkbg}
\color{darkfg}

\title{Project 6: MapReduce and Network Programming}
\author{arika khor}
\date{April 20, 2026}

\begin{document}
\maketitle
\tableofcontents
\newpage

\section{Project Overview}
This project explores large-scale data processing using the MapReduce paradigm and network communication using TCP and UDP protocols. Problem 1 implements a word analyzer for the Gutenberg corpus, while Problem 2 develops a cloud-based multiplication service with robust error handling.

\section{Problem 1: MapReduce Word Length Analyzer}

\subsection{Description}
The objective was to identify the longest words for each letter of the alphabet (a-z) within the text of "Pride and Prejudice" (\texttt{1342-0.txt}).

\subsection{Implementation Details}
\begin{itemize}
    \item \textbf{Framework:} Utilized \texttt{mrjob} for a single-step MapReduce job.
    \item \textbf{Mapper:} Uses \texttt{re.findall(r'[a-zA-Z]+', line)} to extract tokens.
    \item \textbf{Combiner:} Performs local aggregation to find maximum lengths per letter.
    \item \textbf{Reducer:} Performs global aggregation, collecting all words matching the final maximum length into a sorted list.
\end{itemize}

\subsection{Code Snippets}
\begin{lstlisting}[language=Python, caption=MapReduce Logic]
def mapper(self, _, line):
    words = re.findall(r'[a-zA-Z]+', line)
    for word in words:
        yield word[0].lower(), word.lower()

def combiner(self, letter, words):
    max_len = 0
    longest_words = set()
    for word in words:
        current_len = len(word)
        if current_len > max_len:
            max_len = current_len
            longest_words = {word}
        elif current_len == max_len:
            longest_words.add(word)
    for word in longest_words:
        yield letter, word
\end{lstlisting}

\subsection{Execution Output}
\begin{lstlisting}
"g"     ["gentlemanlike", "gratification"]
"h"     ["healthfulness", "hertfordshire"]
"r"     ["recommendations"]
"s"     ["superciliousness"]
"t"     ["thoughtlessness"]
\end{lstlisting}

\section{Problem 2: Network Programming - Multiplication Servers}

\subsection{Description}
This problem implements a "Cloud Multiplication" service where clients send comma-separated numbers to a server and receive their product.

\subsection{Implementation Details}
\begin{itemize}
    \item \textbf{Server Architecture:} Built using \texttt{socketserver.ThreadingTCPServer} and \texttt{ThreadingUDPServer}.
    \item \textbf{Robustness:} Enabled allow\_reuse\_address
    \item \textbf{Multiplication Logic:} Uses \texttt{math.prod()} on a list of floats.
\end{itemize}

\subsection{Code Snippets}
\begin{lstlisting}[language=Python, caption=Multiplication Logic]
def calculate_product(data_str):
    try:
        tokens = [t.strip() for t in data_str.split(',') if t.strip()]
        if not tokens: return "Invalid input"
        numbers = [float(t) for t in tokens]
        product = math.prod(numbers)
        return str(int(product)) if product == int(product) else str(product)
    except (ValueError, TypeError):
        return "Invalid input"
\end{lstlisting}

\subsection{Verification Results}
\begin{table}[h]
\centering
\color{darkfg}
\caption{Client-Server Verification}
\vspace{0.1in}
\begin{tabular}{lll}
\toprule
Input String & Protocol & Server Response \\
\midrule
1, 2, 3 & TCP & 6 \\
1.5, 2 & TCP & 3 \\
hello & TCP & Invalid input \\
10, 20 & UDP & 200 \\
0.5, 10 & UDP & 5 \\
1, a, 2 & UDP & Invalid input \\
\bottomrule
\end{tabular}
\end{table}

\section{Security Considerations \& Future Enhancements}
\begin{itemize}
    \item \textbf{TLS/SSL:} Encrypting the TCP stream.
    \item \textbf{Handshake Protocols:} Implementing a PAKE like \textbf{SPAKE2}.
    \item \textbf{DTLS:} Securing the UDP datagrams.
\end{itemize}

\section{Problem 3: Problem Feedback}
\subsection{Complaints:}
\begin{itemize}
    \item Autograder: Why did we have to setup the file structure from scratch \textit{every} time? It was annoying and didn't really serve much of a purpose. It would've been nice to just clone the entire setup and start coding in the files that the autograder wanted to see instead of having to mess with my file structure to get the autograder to work. This either could have been fixed by a simple git clone or actual clone on repo side, I'm just surprised that we had to clone and do a lot of busy work for every project (unless I missed out on something).
    \item Python venv: Dealing with python venv on OSCAR was fairly annoying. It would have been nice if we could have cloned a pre-set venv, but understandable if it is like something you want the students to figure out on their own. 
    \item Bug reporting: Feels a little dumb? If there was a major bug we would only know via word of mouth. We had to change it ourselves (ok not that bad but I feel like it could have been optimized better student side). If we cloned from a git repo it would've been a lot easier to retroactively update the repo acording to the bugfixes from each project instead of each student needing to update their own version (from a git pull). 
    \item I feel like major bugs should have been announced to the class instead of us having to figure things out by looking in discord threads or announcement threads.
\end{itemize}

\subsection{General thoughts: }
\begin{itemize}
    \item Projects personally felt like they ramped in difficulty up until Project 6. However, I was still able to learn a fair bit! With the increase in difficulty came with more freedom to implement the projects in different ways. They became a fair bit more enjoyable.  
    \item The jump in learning concepts from Project 3 to 4 felt quite large, although it is likely more due to the fact that prior knowledge could get me through Projects 1-3. Still, Project 4 was very interresting to implment!
    \item I wish Project 5 went even more into parallel sorting algorithms as I found that very interesting and felt like it could be applied in a lot of places. For example it would've been nice to look at parallelized quick sort, or took a look at parallelized square root (though that's quite niche)
\end{itemize}


\section{Problem 4: Lecture Feedback}
I did a combination of a(attend in person), and d (watch slide narration recordings async). The slides helped a fair bit when working on the projects. Went with this option given that I could schedule additional work shifts.

\begin{itemize}
    \item I like the weekly quizzes! Good to learn and prepare for the final. Sometimes they felt a little vague in wording but I did like that we had the review portion after each quiz. 
    \item I would've taken an MCQ final on paper rather than being forced to install respondus lockdown browser as I don't want to install malware and its a massive pain to have to dual boot every time I just want to take an exam since it doesn't work on most Linux distros...
    \item It would've been nice to have a 10\% midterm just so we can have a feel for what the final would look like instead of just dropping a final on us last minute? Granted I'm fine with the actual quiz -> MCQ format final, just the way it was announced and made aware felt a little off.
    
\end{itemize}


\end{document}
